\section{An explicit Weil class: the exact minimum at every good prime}
\label{sec:local-weil}

This section fixes the data that the bare phrase ``the Weil class of T5'' leaves unspecified.  Put
\[
 K=\mathbb Q(\omega),\qquad \omega^2+\omega+1=0,\qquad
 E(\mathbb C)=\mathbb C/(\mathbb Z+\mathbb Z\omega),\qquad
 B=E^2,\quad A=B\times B.
\]
Let $\theta_i$ be the degree-one principal polarization on the $i$th elliptic factor, let $\lambda=\theta_1+\theta_2$ on $B$, and let $\theta=\sum_{i=1}^4\theta_i$ on $A$.  The specified Weil action is
\[
 \eta(k)=\operatorname{diag}(k,k,\bar k,\bar k).
\]
Thus its signature is $(2,2)$.  Degrees in this section mean $\deg_\theta V=\int_V\theta^2$ for a surface; no factorial is divided out.  For a class $t$ let
\[
 c_\theta(t)=\min\{\max(N,\deg_\theta Z^+,\deg_\theta Z^-):
 N\in\mathbb Z_{>0},\ Z^\pm\text{ effective integral cycles},\
 [Z^+]-[Z^-]=Nt\}.
\]
Equality is in rational cohomology.  In a good special fiber it can be required in one fixed prime-to-characteristic realization or simultaneously in all realizations; our exact example has the same answer under either convention.

\subsection{The projector and explicit bounded representatives}

\begin{proposition}[P06: an integral multiple of the Weil projector]
\label{prop:P06}
On $H^4(A,\mathbb Q)$ put $R=\eta(\omega)^*$ and $S=\eta(\sqrt{-3})^*$.  The projector onto the rational Weil plane
$W=\bigwedge_K^4 H^1(A,\mathbb Q)$ is
\[
 e_W=\frac{(9+S)(2-R-R^2)}{54}.
\]
The same formula defines its specialized two-dimensional subspace in every good fiber.
\end{proposition}
\begin{proof}
After splitting $K$, write $H^1=V_\sigma\oplus V_{\bar\sigma}$, with both summands of dimension four.  On
$\bigwedge^aV_\sigma\otimes\bigwedge^{4-a}V_{\bar\sigma}$ the eigenvalues of $S$ and $R$ are respectively $9(-1)^a$ and $\omega^{2-a}$.  The first factor $(9+S)/18$ kills odd $a$.  On the remaining values $a=0,2,4$, the second factor $(2-R-R^2)/3$ has eigenvalues $1,0,1$.  Its image is precisely the two Weil lines.  The calculation is in rational endomorphism correspondences and therefore also applies after specialization, using a prime-to-characteristic coefficient field.
\end{proof}

For norm-one $a,b\in\mathbb Z[\omega]$ write
\[
 Z_{a,b}=\{(x,y,ax,by):x,y\in E\},\qquad \Gamma_a=Z_{a,a}.
\]
Each such surface has degree eight.  Put $w=e_W[\Gamma_1]$.  For a pair of factors $i,j$, let $D_{ij}(a)$ be the cross K\"unneth component of the graph divisor $x_j=ax_i$; for a unit it is the graph class minus $\theta_i+\theta_j$.  Set
\[
 P=D_{13}(1),\quad Q=D_{13}(\omega),\quad U=D_{24}(1),\quad V=D_{24}(\omega).
\]
The map $a\mapsto D_{ij}(a)$ is additive, since its cross component is the action of $a$ on $H^1$.

\begin{proposition}[P07: an explicit six-surface certificate]
\label{prop:P07}
One has
\begin{align*}
 3w&=PU-PV-QU-2QV\\
 &= [Z_{1,-\omega}]+[Z_{1,-\omega^2}]+[Z_{\omega,-1}]
 -[Z_{-\omega,-\omega}]-[Z_{-\omega^2,-1}]-[Z_{-\omega^2,-\omega^2}].
\end{align*}
Consequently $c_\theta(w)\leq24$.  More generally, for integers $a,b,m$ with $m>0$, the class $(aw+bRw)/m$ has a representative with multiplier $3m$ and each numerator degree at most $24(|a|+|b|)$.  These certificates specialize with the same bounds at every good place.
\end{proposition}
\begin{proof}
The projector kills the non-cross even terms and the odd terms in $[\Gamma_1]$, leaving $(2-R-R^2)PU/3$.  The identities
$RP=-P-Q$, $RQ=P$, $RU=-U-V$, and $RV=U$ give the first formula.  Expand
\[
 [Z_{a,b}]=(\theta_1+\theta_3+D_{13}(a))
            (\theta_2+\theta_4+D_{24}(b))
\]
for the six indicated surfaces and use $D(\omega^2)=-D(1)-D(\omega)$ and $D(-a)=-D(a)$.  Constant and linear terms cancel; the remaining expression is $PU-PV-QU-2QV$.  Each side contains three degree-eight surfaces.  Pullback by the polarization-preserving automorphism defining $R$ preserves this bound; taking signed integral linear combinations proves the general assertion.  All identities arise from divisor classes and endomorphisms defined over a number field, and hence commute with good specialization.
\end{proof}

Our exact example is
\begin{equation}
\label{eq:tstar}
 t_*=[\Gamma_1]+[\Gamma_{-1}]-[\Gamma_\omega]-[\Gamma_{-\omega}]
       =2(PU-QV)=4w+2Rw.
\end{equation}
This is an explicitly normalized class, not a claim that T5 supplied a particular normalization.  The lattice calculation gives
\[
 w^2=2,\qquad w\cdot Rw=-1,\qquad t_*^2=24.
\]
For completeness, $P^2=Q^2=-2\theta_1\theta_3$, $PQ=\theta_1\theta_3$, and the analogous formulas for $U,V$ give $(PU-QV)^2=6$, proving the last equality directly.  In the ordered integral one-form basis
$(du_1,dv_1,\ldots,du_4,dv_4)$, the coefficient of
$du_1\wedge du_2\wedge dv_3\wedge dv_4$ in $PU-QV$ is $-1$.
Thus $t_*/2$ is primitive integral, while $t_*$ is twice a primitive integral class.  The exact-rational exterior-algebra certificate accompanying this paper checks all 70 coordinates, the projector, the six-surface identity, and these intersection numbers.

\subsection{The exact complex minimum}

Choose holomorphic one-forms $\alpha_i$ with
$(i/2)\int_E\alpha_i\wedge\bar\alpha_i=1$, and put
\[
 \Psi=\alpha_1\wedge\alpha_2\wedge\bar\alpha_3\wedge\bar\alpha_4,
 \qquad z(V)=\frac14\int_V\Psi.
\]
Integrals on singular surfaces are taken on a resolution.  Since
$\int_{E^2}\alpha_1\wedge\alpha_2\wedge\bar\alpha_1\wedge\bar\alpha_2=4$,
\begin{equation}
\label{eq:z-normalization}
 z(\Gamma_a)=\bar a^2,\qquad z(t_*)=2(1-\omega),\qquad |z(t_*)|=2\sqrt3.
\end{equation}

\begin{lemma}[P08: a degree-period estimate over $\mathbb C$]
\label{lem:P08}
For an integral surface $V\subset B\times B$, let $a,b$ be its two projection degrees, with value zero for a nondominant projection.  On a resolution let $c$ be the intersection of the two pulled-back polarizations.  Then
\[
 |z(V)|\leq\sqrt{ab},\qquad c\geq2\sqrt{ab},\qquad
 \deg_\theta V=2(a+b+c).
\]
In particular every effective integral surface cycle has even degree, and
\begin{equation}
\label{eq:small-period-C}
 \deg_\theta T\leq14\quad\Longrightarrow\quad |z(T)|\leq\sqrt2
\end{equation}
for every effective cycle $T$.
\end{lemma}
\begin{proof}
The two pulled-back holomorphic two-forms have positive Hermitian norms squared $4a$ and $4b$.  Cauchy--Schwarz proves the first inequality.  The pulled-back nef divisor classes have squares $2a$ and $2b$; the surface Hodge index theorem gives $c^2\geq4ab$, while $c\geq0$.  Expansion of their sum proves the degree formula.  The nef version follows from the ample version by adding $\varepsilon H$ and taking a limit; the classical input is the surface Hodge index theorem \cite[Chapter V, Theorem 1.9]{Hartshorne}.

A component with nonzero period has $a,b\geq1$ and therefore degree at least eight.  In a cycle of degree at most fourteen there is at most one such component, of multiplicity one.  For that component $a+b+c\leq7$.  If $ab\geq3$, then $a+b\geq4$ and the integral number $c$ is at least four, a contradiction.  Thus $ab\leq2$, proving \eqref{eq:small-period-C}.  This argument includes arbitrary singular components and all positive integral multiplicities.
\end{proof}

\begin{theorem}[P09: the exact minimum over $\mathbb C$]
\label{thm:P09}
The class $t_*$ is a nonzero Weil class and $c_\theta(t_*)=16$.
\end{theorem}
\begin{proof}
Adding the graph correspondences for $a$ and $-a$ cancels odd K\"unneth components.  Subtracting the sums for $a=1$ and $a=\omega$ cancels the degree-zero and degree-four components, since all maps are automorphisms.  On $H^2(B)$ scalar $\omega$ acts as $\omega^2,1,\bar\omega^2$ on $H^{2,0},H^{1,1},H^{0,2}$ respectively.  Only
\[
 H^{2,0}(B)\otimes H^{0,2}(B)
 \ \oplus\ H^{0,2}(B)\otimes H^{2,0}(B)
\]
remains.  For the specified mixed $K$ action these are the two Weil lines.  Nonvanishing follows from \eqref{eq:z-normalization}.

The four graphs in \eqref{eq:tstar} have degree eight, giving $N=1$ and both numerator degrees sixteen.  Conversely, if $Nt_*=[T^+]-[T^-]$ with $N\geq1$ and both effective degrees at most fourteen, then
\[
 2N\sqrt3=|z(T^+)-z(T^-)|\leq2\sqrt2,
\]
which is impossible.  All degrees are even, so a degree bound fifteen is the same as fourteen.  Every multiplier is excluded below sixteen, proving the assertion.
\end{proof}

\subsection{An algebraic replacement for the period at finite primes}

The lower bound just proved cannot be transported by assuming that holomorphic-form positivity survives unchanged in finite characteristic.  We instead use an algebraic coordinate and classify the low-degree surfaces.  This also determines the minimum at the individual small primes.

Let $E$ now denote any geometric good reduction of the specified CM elliptic curve, with its $\mathbb Z[\omega]$ action.  The projector of Proposition~\ref{prop:P06} and the classes $w,Rw$ define a two-dimensional rational subspace.  For any algebraic cycle $Y$, the coefficients of $e_W[Y]$ in this basis are rational: pair with $w,Rw$ and invert their nonsingular rational intersection matrix.  All the pairings are algebraic intersection numbers.  Thus $z$ continues algebraically to a $K$-valued functional, independent of the realization and agreeing with \eqref{eq:z-normalization}.

In the supersingular case write
\[
 D=\operatorname{End}^0(E)=K\oplus Kj,\qquad
 j^2=-\delta,\quad j\alpha=\bar\alpha j,\quad \delta>0.
\]
The reduced norm is $\mathrm N(\alpha+\beta j)=|\alpha|^2+\delta|\beta|^2$; write $q_K=\alpha$ and $q_{\rm anti}=\beta j$.  The ordinary case has $D=K$.  These are the standard endomorphism-algebra alternatives for elliptic curves \cite[Theorem II.6.9]{MilneEC}.

Two formulas will be used.  For a homomorphism $h:E^2\to E^2$,
\begin{equation}
\label{eq:graph-coordinate}
 z(\Gamma_h)=\overline{\det(h_K)}.
\end{equation}
For graph curves $C_q,C_r$ in the two copies of $E^2$, with anti components $\beta j,\gamma j$,
\begin{equation}
\label{eq:product-coordinate}
 z(C_q\times C_r)=-\delta\beta\bar\gamma,
 \qquad |z(C_q\times C_r)|^2=\mathrm N(q_{\rm anti})\mathrm N(r_{\rm anti}).
\end{equation}
The sign or conjugate changes with simultaneous conventions and is immaterial to the inequalities; the displayed convention agrees with \eqref{eq:z-normalization}.
To check the constants, choose $K$ eigenvectors $e,f$ of $H^1(E)$ with
$j^*e=f$, $j^*f=-\delta e$, and put $\tau=\operatorname{Tr}_E(e\wedge f)$.
The coordinate is the trace against $e_1e_2f_3f_4$ divided by $-\tau^2$.
For $q=\alpha+\beta j$ one has
$q^*e=\alpha e+\beta f$ and $q^*f=\bar\alpha f-\delta\bar\beta e$.
Expanding proves \eqref{eq:graph-coordinate}; for the two graph curves the traces are $\beta\tau$ and $\delta\bar\gamma\tau$, proving \eqref{eq:product-coordinate}.  These identities extend linearly to divisor classes.  Translations act trivially, and interchanging the two surface factors conjugates $z$.

\begin{lemma}[Low-degree geometric reductions used in P13]
\label{lem:lowdegree-geometry}
Let $V\subset E^2\times E^2$ be an integral surface of degree at most fourteen.
If both projections dominate, then $V$ is a translate of a graph or transpose graph of an isogeny $h:E^2\to E^2$ of degree $d\in\{1,2\}$, and
\begin{equation}
\label{eq:graph-degree-budget}
 \deg_\theta\Gamma_h=2(1+d+T_h),\qquad
 T_h=\sum_{i,j}\mathrm N(h_{ij})\leq6-d.
\end{equation}
If a projection has image a curve $C$, its middle K\"unneth component is $[C]\otimes[F]$, where $F$ is an effective general-fiber divisor class in the other $E^2$, and
\begin{equation}
\label{eq:fiber-degree-budget}
 \deg_\theta V\geq2\deg_\lambda(C)\deg_\lambda(F).
\end{equation}
A point image contributes zero to $z$.
\end{lemma}
\begin{proof}
Resolution of surfaces and the surface Hodge index theorem apply in every characteristic.  Their numerical part gives
$\deg_\theta V=2(a+b+c)$ and $c\geq2\sqrt{ab}$ exactly as above.  If $a,b>0$, the same integer threshold gives $ab\leq2$.  A degree-one projection makes the other projection a rational map from $E^2$ to itself.  It extends to a morphism, which is a homomorphism followed by translation; see \cite[Corollary 2.2 and Theorem 3.1]{MilneAV}.  Its closed graph is $V$.  Intersecting $h^*\lambda$ with the two coordinate elliptic curves gives $\lambda h^*\lambda=\sum\mathrm N(h_{ij})$, proving \eqref{eq:graph-degree-budget}.  Degrees include inseparable degrees.

For a curve image, normalize $C$ and take the closure of $V$ over its smooth open part in $\widetilde C\times E^2$.  This is an integral Cartier divisor in a smooth threefold, flat over $\widetilde C$, and maps birationally onto $V$.  Its fiber divisor has constant class $[F]$.  In the K\"unneth decomposition of its first Chern class, the term $1\otimes[F]$ pushes to $[C]\otimes[F]$; the other terms push into bidegrees $(3,1)$ and $(4,0)$ and do not contribute to $W$.  The mixed intersection equals $\deg_\lambda(C)\deg_\lambda(F)$, proving the inequality.  This construction introduces no generic multiplicity.
\end{proof}

We also need two elementary consequences.  An effective divisor of degree one in $E^2$ is a coordinate fiber.  An integral divisor of degree two with both elliptic projections nonconstant has projection degrees $(1,1)$ and is a translated automorphism graph.  An integral divisor of degree three has degrees $(1,2)$ or $(2,1)$ and is a graph or transpose graph of a degree-two endomorphism.  If one elliptic projection is constant, an integral divisor is a coordinate fiber.  This classifies effective divisor cycles of degree at most three by their components and multiplicities.

Second, if a row or column of a quaternionic matrix $h$ lies in $K$, then
\begin{equation}
\label{eq:schur-bound}
 |\det(h_K)|^2\leq\deg h.
\end{equation}
Indeed Schur factorization at a nonzero $K$ pivot expresses the degree as the norm of that pivot times the norm of the quaternionic Schur complement.  Projection to $K$ is $K$-bilinear and decreases norm, giving the ordinary determinant of $h_K$.  A zero pivot is handled by a permutation.  If any entry of $h$ is zero, a row and column permutation makes it triangular and the same inequality follows from the product of the two diagonal norms.  The identification of degree with the determinant on the Tate module is \cite[Proposition 12.9]{MilneAV}.

\begin{theorem}[P13: exact complexity at every geometric good reduction]
\label{thm:P13}
For the specified product polarization and class $t_*$,
\[
                         c_\theta(t_{*,v})=16
\]
at every geometric good reduction, ordinary or supersingular.  This includes residue characteristics two and, when a good model is taken after a field extension, three.
\end{theorem}
\begin{proof}
The graph expression gives the upper bound in every characteristic.  We prove that every effective cycle of degree at most fourteen has $|z|<\sqrt3$.

\emph{Residue characteristic $p>3$.}
For $j(E)=0$, every automorphism is in $\mu_6\subset K$ \cite[Proposition IV.7.13]{MilneEC}.  A non-$K$ integral endomorphism consequently has integral norm at least two.  In \eqref{eq:graph-degree-budget}, three such entries would cost at least six.  If exactly two occupy opposite corners, at most one further entry can be nonzero, forcing determinant degree at least four.  If they occupy one row or column, at budget five the only nonsingular possibility has one additional unit and has degree at least two, contradicting $d=1$; at budget four the remaining row or column is zero, contradicting $d=2$.  Thus a low-degree graph has at most one non-$K$ entry, and \eqref{eq:schur-bound} gives $|z|\leq\sqrt2$.  In the ordinary case all entries already lie in $K$.

For a nondominant projection, a degree-one or degree-two divisor has zero anti-$K$ component by the preceding curve classification and the automorphism description.  If $z(V)\ne0$, both divisors in \eqref{eq:fiber-degree-budget} have degree at least three.  The surface degree is then at least eighteen.  Thus all nondominant surfaces of degree at most fourteen have zero $z$.

\emph{Residue characteristic two.}
The endomorphism order is the Hurwitz order
\[
 \mathcal O_H=\{q_0+q_1i+q_2j+q_3k:
  q_i\text{ all integers or all half-integers}\},
 \qquad i^2=j^2=k^2=ijk=-1.
\]
For this classical identification see \cite[footnote 2]{CHVW2022}.
Choose $\omega=(-1+i+j+k)/2$.  Any embedding of the Eisenstein order is carried to this one by an orientation-preserving signed permutation of the three imaginary axes, which preserves the Hurwitz order: its trace-minus-one norm-one generator has exactly these half-integral coordinates with arbitrary imaginary signs.  The $K$ projection is
\[
 q_K=q_0+\frac{q_1+q_2+q_3}{3}(i+j+k).
\]
The 24 units are $\pm1,\pm i,\pm j,\pm k$ and $(\pm1\pm i\pm j\pm k)/2$.  Their $K$-projection moduli are one for the six $K$ units and $1/\sqrt3$ for the others.  Norm-two elements have exactly two integral coordinates $\pm1$, and therefore $K$-projection modulus at most $2/\sqrt3$; none lies in $K$.

A zero-entry graph matrix is covered by \eqref{eq:schur-bound}.  Otherwise the allowed norm patterns are $(1,1,1,1)$ and $(2,1,1,1)$.  For four units, all entries in $K$ give $|z|\leq\sqrt2$; otherwise the determinant triangle inequality gives $|z|\leq1+1/\sqrt3<5/3$.
The second pattern occurs only for $d=1$.  A full $K$ row or column gives $|z|\leq1$ by \eqref{eq:schur-bound}.  If neither exists, put the norm-two entry at the upper left.  If the opposite unit is non-$K$, the determinant products have moduli at most $2/3$ and one, giving $|z|\leq5/3$.  If the opposite unit is in $K$, both adjacent units must be non-$K$, giving $|z|\leq2/\sqrt3+1/3<5/3$.

For a nondominant surface with nonzero $z$, the two divisor degrees are $(2,2),(2,3)$, or $(3,2)$ by \eqref{eq:fiber-degree-budget}.  A degree-two non-$K$ unit graph has anti-$K$ norm $1-1/3=2/3$.  A degree-three divisor has anti-$K$ norm at most two: it is a norm-two graph, a unit graph plus a coordinate fiber, or a sum of coordinate fibers.  Formula~\eqref{eq:product-coordinate} gives $|z|\leq\sqrt{(2/3)2}=2/\sqrt3<5/3$.

\emph{Residue characteristic three.}
Over an algebraic closure take $E:y^2=x^3-x$.  The automorphisms
\[
 \omega:(x,y)\longmapsto(x+1,y),\qquad
 i:(x,y)\longmapsto(-x,s y),\quad s^2=-1,
\]
satisfy $i^2=-1$ and $i\omega=\bar\omega i$.  The order
$\mathbb Z[\omega]\oplus\mathbb Z[\omega]i$ has reduced discriminant three: the absolute determinant of its trace Gram matrix on $1,\omega,i,\omega i$ is nine.  It is maximal in the quaternion algebra ramified at three and infinity and hence equals the endomorphism order.  Its norm is
$\mathrm N(A+Bi)=\mathrm N_K(A)+\mathrm N_K(B)$.
Norm-one elements have $K$-projection modulus one or zero; norm-two elements have projection modulus one, since there is no Eisenstein integer of norm two.
For the four-unit pattern, either all entries are in $K$ or a determinant product vanishes, giving $|z|\leq\sqrt2$ or $|z|\leq1$.  In the $(2,1,1,1)$ pattern, either all three units are in $K$ and \eqref{eq:schur-bound} applies, or one projection vanishes; again $|z|\leq1$.  Nondominant surfaces have divisor degrees at most three as above, with anti-$K$ norm bounds one and two, and therefore $|z|\leq\sqrt2$.

\emph{Effective sums and all multipliers.}
In all cases a $z$-nonzero integral component has degree at least eight: use the two-dominant-projection intersection inequality or the product of two divisor degrees at least two.  An effective cycle of degree at most fourteen therefore has at most one such component, with multiplicity one.  The bounds just proved are $\sqrt2$ for $p\ne2$ and $5/3$ for $p=2$, both strictly smaller than $\sqrt3$.  All surface degrees are even by the numerical degree formula.  If $Nt_{*,v}=[T^+]-[T^-]$ with $N\geq1$ and numerator degrees below sixteen, then
\[
 2N\sqrt3=|z(T^+)-z(T^-)|<2\sqrt3,
\]
a contradiction.  The upper graph certificate attains sixteen.  This argument uses only algebraic intersection numbers and a prime-to-characteristic realization; its upper identity comes from global cycles and holds simultaneously in all realizations.
\end{proof}

\begin{proposition}[P14: the unrestricted positive-characteristic inequality is false]
\label{prop:P14}
The inequality $|z(V)|\leq\sqrt{ab}$ of Lemma~\ref{lem:P08} does not hold for all supersingular surfaces.
\end{proposition}
\begin{proof}
At $p=5$ let $j$ be Frobenius of the trace-zero $j=0$ elliptic curve.  Then $j^2=-5$ and $j\omega=\bar\omega j$.  The matrix
\[
 h=\begin{pmatrix}2&j\\-j&2\end{pmatrix}
\]
is an automorphism: its product with $\left(\begin{smallmatrix}2&-j\\j&2\end{smallmatrix}\right)$ is $-I$.  Its graph has both projection degrees one, while $h_K=2I$ gives $z(\Gamma_h)=4$.  Its degree is $2(1+1+4+5+5+4)=40$.  Thus it refutes the unrestricted inequality, and leaves the low-degree argument of Theorem~\ref{thm:P13} intact.
\end{proof}

\subsection{Normalization, prime types, and the family distinction}

\begin{corollary}[P12: exact polarization dependence]
\label{cor:P12}
For every integer $k\geq1$,
\[
 c_{\theta^{\otimes k}}(t_*)=16k^2
\]
over $\mathbb C$ and at every geometric good reduction.  Over $\mathbb C$, for every integer $M\geq1$,
\[
 8\sqrt3\,M\leq c_\theta(Mt_*)\leq16M.
\]
\end{corollary}
\begin{proof}
Surface degrees are multiplied by $k^2$.  The proofs above excluded every positive multiplier with both original degrees below sixteen, whereas the multiplier-one graph certificate attains sixteen.  For the second assertion, Lemma~\ref{lem:P08} gives $|z(T)|\leq\deg_\theta(T)/8$ by $a+b\geq2\sqrt{ab}$ and $c\geq2\sqrt{ab}$, followed by additivity.  A difference representing $NMt_*$ therefore requires $D\geq8\sqrt3 NM$.  The upper graph certificate is multiplied by $M$.
\end{proof}

This is a concrete reason that a least complexity cannot be a function of the prime and Newton type alone when the polarization or class normalization is left variable.  The exact answer sixteen concerns \eqref{eq:tstar}; an arbitrary flat Weil class on an arbitrary T5 fiber is a separate minimization problem.  For example the explicitly normalized $w$ above satisfies $16\leq c_\theta(w)\leq24$ over $\mathbb C$: one has $z(w)=1$, and its denominator in integral cohomology is exactly three.  Thus every allowable multiplier is at least three, and the lower argument in Lemma~\ref{lem:P08} excludes degree at most fourteen since $3>2\sqrt2$.  Proposition~\ref{prop:P07} supplies the upper bound.  The certificate verifies primitiveness of $3w$; one coefficient of $3w$ in the indicated exterior basis is $-1$.

\begin{proposition}[P02: prime types on the fixed $E^4$ isogeny orbit]
\label{prop:P02}
If an abelian fourfold over a number field is geometrically isogenous to $E^4$, then at every common good place of residue characteristic $p>3$ it is ordinary for $p\equiv1\pmod3$ and supersingular for $p\equiv2\pmod3$.
\end{proposition}
\begin{proof}
Use the $j=0$ model $y^2=x^3+1$.  The Hasse-invariant criterion \cite[Chapter V, Theorem 4.1(a)]{SilvermanAEC} tests the coefficient of $x^{p-1}$ in $(x^3+1)^{(p-1)/2}$.  It vanishes for $p\equiv2\pmod3$; for $p\equiv1\pmod3$ it is
$\binom{(p-1)/2}{(p-1)/3}$, nonzero modulo $p$.  These are exactly the supersingular and ordinary elliptic cases.  Powers repeat the slopes.  An isogeny, after a finite extension where it is defined, specializes between good abelian schemes and preserves the rational Dieudonn\'e module and Newton polygon.  Residue-field extensions do not change these slopes.
\end{proof}

The two infinite subsequences therefore have the identical least complexity $16$ for our fixed class.  Using the Dirichlet GRH theorem cited in the prompt, their rational-prime counts are
\[
 \pi(x;3,a)=\tfrac12\operatorname{Li}(x)+O(\sqrt{x}\log x),\qquad a=1,2,
\]
after omission of finitely many bad primes.  An exact primary locator for this standard consequence is \cite[Corollary 1.1, equation (2)]{Lee2023}; only its fixed-modulus asymptotic consequence is used.  Infinitude needs no GRH.  Counting number-field places by their norms is a different assertion and is not silently identified with rational-prime counting.
The attached \texttt{ALPOGE\_Main.lean} and \texttt{ALPOGE\_Statements.lean} have precisely Dirichlet-character scope.  The general Hecke input remains conditional as in the sketch: the supplied Masullo index at PDF p.~522, Theorem~13.3.1.2, contains an unfilled main-statement entry.  No imported Lean proof was rebuilt in this run.

\begin{proposition}[P03: the whole Weil family has a different prime-type problem]
\label{prop:P03}
Splitting in $K=\mathbb Q(\omega)$ does not determine the Newton type of every abelian fourfold of Weil type for $K$.
\end{proposition}
\begin{proof}
For any elliptic curve $C$ act on each of two $C^2$ blocks by
\[
 M=\begin{pmatrix}0&-1\\1&-1\end{pmatrix},\qquad
 H=\begin{pmatrix}2&-1\\-1&2\end{pmatrix}.
\]
The positive polarization matrix $H$ satisfies $M^{\mathsf t}HM=H$ and $M^2+M+I=0$.  Each complex Lie block has one copy of each embedding of $K$, giving signature $(2,2)$ on $C^4$.  Take $C:y^2=x^3-x$.  It is ordinary at five, since the $x^4$ coefficient of $(x^3-x)^2$ is $-2\not\equiv0\pmod5$, although five is inert in $K$.  It is supersingular at seven, since the $x^6$ coefficient of $(x^3-x)^3$ is zero, although seven splits in $K$.  Both reductions are good.  This disproves the claimed extension of the split/inert rule from the particular isogeny orbit to the entire class of Weil families.
\end{proof}

For the T5 field, the algebraicity input itself predates the cited CM theorem: Schoen's unnumbered theorem on printed p.~329 of the 1998 addendum covers every Eisenstein Weil fourfold \cite{Schoen1998}; that addendum corrects the polarization qualification in the 1988 Theorem~3.2.  Markman's Theorem~1.5.1 and Corollary~1.6.1 give the broader sixfold and fourfold conclusions \cite{Markman2025}.  These source theorems give fixed-class upper bounds by specialization and family bounds after the bounded-Chow argument.  Prime counting alone supplies no degree parameter: the additional geometric statement is that one fixed bounded-degree locus contains the relevant representatives.  The explicit local result above makes that distinction testable without claiming a solution of the higher-dimensional bounded-division problem.
